\documentclass[11pt]{article}
\usepackage{amsmath,amssymb}
\usepackage[a4paper,margin=1in]{geometry}
\usepackage{hyperref}

% Remove paragraph indentation
\setlength{\parindent}{0pt}
\setlength{\parskip}{6pt}

\begin{document}

\title{A Concise Derivation of Linear Discriminant Analysis}
\author{Zhenwei Tang}
\date{}
\maketitle

% -------------------------------------------------------
\section*{Goal}
% -------------------------------------------------------

Given $n$ labeled data points
\[
  \{(x_i, c_i)\}_{i=1}^{n}, \quad x_i \in \mathbb{R}^p, \quad c_i \in \{1, \dots, K\},
\]
where $K$ is the number of classes, LDA seeks a linear projection
\[
  y = w^T x, \quad w \in \mathbb{R}^p, \quad y \in \mathbb{R},
\]
such that the projected data are \textbf{maximally separated between classes} while being \textbf{maximally compact within each class}.
This single criterion is captured by the \emph{Fisher ratio}:
\[
  J(w) = \frac{w^T S_B\, w}{w^T S_W\, w},
\]
where $S_B$ is the \emph{between-class scatter matrix} and $S_W$ is the \emph{within-class scatter matrix} (both defined below).
The derivation proceeds in four steps: define the scatter matrices, formulate the optimization, apply the generalized eigenvalue technique, and read off the solution.

% -------------------------------------------------------
\section*{Step 1: Define the Scatter Matrices}
% -------------------------------------------------------

\textbf{Notation.}
\begin{itemize}
  \item $n_k$ = number of samples in class $k$.
  \item $\mu_k = \dfrac{1}{n_k}\displaystyle\sum_{i:\,c_i=k} x_i$ \quad (class mean).
  \item $\mu = \dfrac{1}{n}\displaystyle\sum_{i=1}^n x_i = \dfrac{1}{n}\sum_{k=1}^K n_k\,\mu_k$ \quad (global mean).
\end{itemize}

\textbf{Within-class scatter matrix} $S_W \in \mathbb{R}^{p\times p}$: measures the total spread of samples around their own class mean.
\[
  S_W = \sum_{k=1}^{K} \sum_{i:\,c_i=k} (x_i - \mu_k)(x_i - \mu_k)^T.
\]

\textbf{Between-class scatter matrix} $S_B \in \mathbb{R}^{p\times p}$: measures how far the class means are from the global mean, weighted by class size.
\[
  S_B = \sum_{k=1}^{K} n_k\,(\mu_k - \mu)(\mu_k - \mu)^T.
\]

Both $S_W$ and $S_B$ are symmetric positive semi-definite (see Appendix~\ref{app:psd}).

After projection onto direction $w$, the projected class mean is $w^T\mu_k$ and the projected global mean is $w^T\mu$.
The scalar between-class and within-class scatters become, respectively,
\[
  \widetilde{s}_B = \sum_{k=1}^K n_k\,(w^T\mu_k - w^T\mu)^2 = w^T S_B\, w,
\]
\[
  \widetilde{s}_W = \sum_{k=1}^K\sum_{i:\,c_i=k}(w^T x_i - w^T\mu_k)^2 = w^T S_W\, w.
\]

% -------------------------------------------------------
\section*{Step 2: Formulate the Optimization Problem}
% -------------------------------------------------------

We wish to find $w$ that maximizes class separation relative to within-class spread.
Formally:
\[
  \max_{w \in \mathbb{R}^p} \; J(w) = \frac{w^T S_B\, w}{w^T S_W\, w}.
\]

\textbf{Why the ratio?}
Maximizing only $w^T S_B w$ is meaningless because we can make it arbitrarily large by scaling $w$.
Dividing by $w^T S_W w$ removes the scale ambiguity and measures \emph{relative} separation.

\textbf{Scale invariance.}
$J(\alpha w) = J(w)$ for any $\alpha \neq 0$, so we may impose the normalization
\[
  w^T S_W\, w = 1
\]
without loss of generality.
The problem becomes an equivalent constrained optimization:
\[
  \max_{w}\; w^T S_B\, w \quad \text{subject to} \quad w^T S_W\, w = 1.
\]

% -------------------------------------------------------
\section*{Step 3: Apply the Lagrange Multiplier Method}
% -------------------------------------------------------

\textbf{Using the Lagrange multiplier method} (see Appendix~\ref{app:lagrange}), we introduce $\lambda \in \mathbb{R}$ and form the Lagrangian:
\[
  \mathcal{L}(w, \lambda) = w^T S_B\, w - \lambda\,(w^T S_W\, w - 1).
\]

\textbf{Taking the derivative with respect to $w$ and setting it to zero} (see Appendix~\ref{app:matderiv}):
\[
  \frac{\partial \mathcal{L}}{\partial w} = 2S_B\, w - 2\lambda\, S_W\, w = 0.
\]

This gives the \textbf{generalized eigenvalue equation}:
\[
  \boxed{S_B\, w = \lambda\, S_W\, w.}
\]

\textbf{Why this is useful.}
The stationary points of $J(w)$ are exactly the generalized eigenvectors of the matrix pair $(S_B, S_W)$.
Substituting a solution $w$ back into $J$:
\[
  J(w) = \frac{w^T S_B\, w}{w^T S_W\, w} = \frac{w^T (\lambda S_W\, w)}{w^T S_W\, w} = \lambda.
\]
Hence, the maximum of $J(w)$ equals the largest generalized eigenvalue $\lambda_{\max}$, achieved by the corresponding generalized eigenvector.

% -------------------------------------------------------
\section*{Step 4: Reduce to a Standard Eigenvalue Problem}
% -------------------------------------------------------

If $S_W$ is invertible (i.e., $n > p$ and no class is degenerate), multiply both sides of the generalized eigenvalue equation by $S_W^{-1}$:
\[
  S_W^{-1} S_B\, w = \lambda\, w.
\]
This is a \textbf{standard eigenvalue problem} for the matrix $S_W^{-1}S_B$.

\textbf{Observation on rank.}
$S_B$ is a sum of $K$ rank-1 matrices $(\mu_k - \mu)(\mu_k-\mu)^T$, so
\[
  \operatorname{rank}(S_B) \leq K - 1.
\]
Therefore $S_W^{-1}S_B$ has at most $K-1$ non-zero eigenvalues, and LDA can produce \textbf{at most $K-1$ discriminant directions}.

% -------------------------------------------------------
\section*{Step 5: Extract the Discriminant Directions}
% -------------------------------------------------------

Let $\lambda_1 \geq \lambda_2 \geq \cdots \geq \lambda_{K-1} > 0$ be the non-zero eigenvalues of $S_W^{-1}S_B$, with corresponding eigenvectors $w_1, w_2, \dots, w_{K-1}$.

To project into a $d$-dimensional subspace ($d \leq K-1$), collect the top-$d$ eigenvectors:
\[
  W_d = [w_1 \mid w_2 \mid \cdots \mid w_d] \in \mathbb{R}^{p \times d}.
\]

The low-dimensional representation of a data point $x$ is:
\[
  y = W_d^T x \in \mathbb{R}^d.
\]

Each eigenvector $w_j$ is a \textbf{linear discriminant}: the direction along which the Fisher ratio is the $j$-th largest.

% -------------------------------------------------------
\section*{Conclusion}
% -------------------------------------------------------

LDA seeks a linear projection that maximizes the Fisher ratio
\[
  J(w) = \frac{w^T S_B\, w}{w^T S_W\, w},
\]
balancing between-class separation $S_B$ against within-class compactness $S_W$.
By imposing $w^T S_W w = 1$ and applying the Lagrange multiplier method, the optimality condition reduces to the generalized eigenvalue problem
\[
  S_B\, w = \lambda\, S_W\, w,
\]
equivalently $S_W^{-1}S_B\, w = \lambda\, w$ when $S_W$ is invertible.
The optimal directions are the eigenvectors corresponding to the $d$ largest eigenvalues.

\begin{center}
  \fbox{\parbox{0.88\textwidth}{\centering
    LDA finds the linear directions that \textbf{maximize between-class scatter relative to within-class scatter},
    obtained as the leading eigenvectors of $S_W^{-1}S_B$.
    At most $K-1$ such directions exist for $K$ classes.
  }}
\end{center}

% -------------------------------------------------------
\appendix
\section*{Appendices}
% -------------------------------------------------------

\subsection*{A\quad Lagrange Multiplier Method}
\label{app:lagrange}

Given a constrained optimization problem
\[
  \max_w\; f(w) \quad \text{subject to} \quad g(w) = 0,
\]
the \emph{Lagrange multiplier method} asserts that at any local optimum $w^*$, there exists a scalar $\lambda$ (the \emph{Lagrange multiplier}) such that
\[
  \nabla_w f(w^*) = \lambda\, \nabla_w g(w^*),
\]
i.e., the gradients of the objective and constraint are parallel.
Equivalently, $w^*$ is a stationary point of the \emph{Lagrangian}
\[
  \mathcal{L}(w,\lambda) = f(w) - \lambda\, g(w).
\]
Setting $\nabla_w \mathcal{L} = 0$ gives the necessary first-order conditions.

\subsection*{B\quad Matrix Derivative of a Quadratic Form}
\label{app:matderiv}

For a symmetric matrix $A \in \mathbb{R}^{p\times p}$ and vector $w \in \mathbb{R}^p$,
\[
  \frac{\partial}{\partial w}(w^T A\, w) = 2A\, w.
\]
\textit{Proof sketch.} Writing $(w^T A w) = \sum_{i,j} A_{ij} w_i w_j$, differentiate with respect to $w_k$:
\[
  \frac{\partial}{\partial w_k}\sum_{i,j} A_{ij} w_i w_j
  = \sum_j A_{kj} w_j + \sum_i A_{ik} w_i = (Aw)_k + (A^T w)_k = 2(Aw)_k,
\]
using symmetry $A = A^T$. Hence $\nabla_w (w^TAw) = 2Aw$.

\subsection*{C\quad Symmetric Positive Semi-Definite Matrices}
\label{app:psd}

A matrix $M \in \mathbb{R}^{p\times p}$ is \emph{symmetric positive semi-definite} (PSD), written $M \succeq 0$, if:
\begin{enumerate}
  \item $M = M^T$ (symmetry), and
  \item $v^T M\, v \geq 0$ for all $v \in \mathbb{R}^p$.
\end{enumerate}

Any matrix of the form $M = \sum_i z_i z_i^T$ (a sum of outer products) is PSD, because
\[
  v^T M\, v = \sum_i v^T z_i z_i^T v = \sum_i (z_i^T v)^2 \geq 0.
\]
Both $S_W$ and $S_B$ have exactly this form, so they are PSD.

If additionally $v^T M v > 0$ for all $v \neq 0$, then $M$ is \emph{positive definite} (PD), written $M \succ 0$, and $M$ is invertible.
$S_W \succ 0$ whenever the within-class data span $\mathbb{R}^p$ (i.e., $n > p$ with full-rank scatter), which is the condition required in Step~4 for $S_W^{-1}$ to exist.

\subsection*{D\quad Generalized Eigenvalue Problem}
\label{app:geneig}

Given two symmetric matrices $A, B \in \mathbb{R}^{p\times p}$ with $B \succ 0$, the \emph{generalized eigenvalue problem} is
\[
  A\, v = \lambda\, B\, v.
\]
A scalar $\lambda$ satisfying this (for some $v \neq 0$) is a \emph{generalized eigenvalue}; the corresponding $v$ is a \emph{generalized eigenvector}.
When $B$ is invertible, this reduces to the standard eigenvalue problem $B^{-1}A\, v = \lambda\, v$.
The generalized eigenvectors diagonalize the ratio $v^T A v / v^T B v$, and the maximum of this ratio equals the largest generalized eigenvalue.

\end{document}
